%==============================================================================
%  CAPÍTULO — EFECTOS LABORALES DE LA APERTURA DEL CONCURSO
%==============================================================================
\chapter{Efectos de la apertura del concurso sobre los contratos de trabajo}
\label{cap:laboral-efectos}

\fichaCapitulo{El concurso como causal de término.}{El artículo 163 bis del Código del
Trabajo, la determinación de la fecha del despido, la suerte de los fueros y la posición
del liquidador como empleador.}

%==============================================================================
\bloqueTeorico
%==============================================================================

\subsection{Dos ordenamientos con lógicas opuestas}

El derecho laboral se construye sobre la estabilidad del vínculo y la protección del
trabajador frente al poder del empleador. El derecho concursal se construye sobre la
liquidación ordenada de un patrimonio insuficiente y la distribución paritaria de la
pérdida. Cuando una empresa con trabajadores entra en liquidación, ambos ordenamientos se
encuentran, y el punto de encuentro es el \artcdt{163 bis} \citeyearpar{codigo-trabajo}.

La solución del legislador es tajante: la resolución de liquidación pone término a los
contratos de trabajo por el solo ministerio de la ley. No hay despido en sentido
estricto, sino terminación legal. Esa calificación tiene consecuencias que recorren todo
este capítulo, y que conviene anticipar:

\begin{enumerate}
  \item No se requiere invocar una causal de las del \artcdt{159}, \artcdt{160} o
  \artcdt{161}: la causal es la resolución misma.
  \item La fecha del término es la de dictación de la resolución de liquidación, no la de
  su notificación ni la de la comunicación al trabajador.
  \item Quien debe cumplir los deberes del empleador ---comunicación, finiquito, pago---
  es el liquidador, con cargo a la masa.
\end{enumerate}

\subsection{Por qué el legislador optó por la terminación automática}

La alternativa ---mantener vigentes los contratos hasta que el liquidador decidiera
individualmente sobre cada uno--- habría producido tres efectos indeseables: la
acumulación de remuneraciones devengadas durante la liquidación como créditos de la masa,
la incertidumbre prolongada del trabajador respecto de su situación, y la imposibilidad
práctica de realizar el activo mientras subsistieran obligaciones laborales corrientes.

La terminación automática corta ese nudo, pero traslada el problema al terreno del pago:
los créditos laborales devengados hasta la fecha del término deben pagarse con
preferencia, y esa preferencia es el objeto del capítulo \ref{cap:creditos-laborales}.

%==============================================================================
\bloqueNormativo
%==============================================================================

\subsection{La causal del artículo 163 bis}

\subsubsection{Supuesto de hecho}

La causal opera con la dictación de la resolución de liquidación respecto del empleador.
No se extiende a la resolución de reorganización: durante la reorganización los contratos
de trabajo subsisten, y su terminación requiere causal propia.

\begin{alerta}[Distinción esencial]
Reorganización y liquidación producen efectos laborales opuestos. La reorganización
mantiene los contratos vigentes; la liquidación los termina de pleno derecho. Confundirlas
en la asesoría genera despidos injustificados o, a la inversa, expectativas de
continuidad que no existen.
\end{alerta}

\subsubsection{Fecha del término}

El \artcdt{163 bis} lo dice con toda precisión: \destacar{«para todos los efectos legales, la
fecha de término del contrato de trabajo será la fecha de dictación de la resolución de
liquidación»}. La precisión importa porque de ella dependen el cómputo de la antigüedad, la
determinación de las remuneraciones devengadas y el corte entre créditos concursales y créditos de
la masa.

\subsubsection{Indemnizaciones que procede pagar}

\begin{description}[leftmargin=0pt,style=nextline,font=\sffamily\bfseries]
  \item[Indemnización sustitutiva del aviso previo (\Nro 2)] Equivale al \destacar{promedio de las
  tres últimas remuneraciones mensuales devengadas}; si hay menos de tres, al promedio de las dos
  últimas y, en su defecto, a la última remuneración mensual devengada. El precepto
  \emph{no} contempla tope en unidades de fomento: el límite de \uf{90} opera en sede de
  \emph{preferencia} ---\artcc{2472} \Nro 5---, no de procedencia del crédito (capítulo
  \ref{cap:creditos-laborales}).

  \item[Indemnización por años de servicio (\Nro 3)] Procede si el contrato estuvo vigente
  \destacar{un año o más}, y equivale a la que correspondería si el contrato terminara por las
  causales del \artcdt{161}, determinada conforme a los incisos primero y segundo del
  \artcdt{163}. Su base de cálculo es la \destacar{última remuneración mensual} en los términos del
  \artcdt{172} ---con tope de \destacar{90 UF}---, y el límite temporal es de \destacar{once años}
  (330 días), salvo los trabajadores con contrato vigente anterior al 14 de agosto de 1981, a
  quienes se reconoce la totalidad de sus años de servicio sin tope. Es \destacar{compatible} con
  la indemnización sustitutiva del \Nro 2.

  \item[Feriado proporcional y feriado pendiente] Se devengan hasta la fecha del término.

  \item[Remuneraciones y cotizaciones adeudadas] Con la preferencia que corresponda según
  el capítulo \ref{cap:creditos-laborales}.
\end{description}

\subsection{Los fueros laborales frente a la liquidación}

Es la materia de mayor conflictividad práctica del capítulo. La pregunta es si el fuero
maternal, el fuero sindical y los demás fueros especiales resisten la terminación del
\artcdt{163 bis}, o si ceden ante ella.

La ley especial \destacar{desplaza las normas generales sobre fuero}: la resolución de
liquidación permite el término de los contratos de los trabajadores amparados por fuero
gremial o de maternidad \destacar{sin necesidad de autorización judicial de desafuero previa}.
La causal opera por el solo ministerio de la ley y no admite excepción por fuero, precisamente
porque no hay un despido decidido por el empleador que la voluntad tutelada por el fuero pueda
neutralizar.

Para mitigar el impacto sobre la trabajadora con fuero maternal del \artcdt{201}, el liquidador
debe pagar con cargo a la masa una \destacar{indemnización equivalente a la última remuneración
mensual devengada por cada uno de los meses que restare de fuero}. El texto añade dos precisiones
de cálculo que la práctica suele pasar por alto: \emph{no} se computan las semanas durante las
cuales la trabajadora tenga derecho a los \destacar{subsidios} de los descansos y permisos del
\artcdt{198}; y esta indemnización es \destacar{compatible con la indemnización por años de
servicio} del \Nro 3, pero \destacar{incompatible con la sustitutiva del aviso previo} del \Nro 2.

\begin{foco}[El fuero cede, pero se indemniza]
La regla práctica es doble: el fuero no impide el término ---no hay que desaforar--- pero
genera una indemnización sustitutiva a cargo de la masa. El liquidador que despide a una
trabajadora aforada sin provisionar y pagar esa indemnización incurre en responsabilidad;
quien exige el desafuero previo, en cambio, paraliza innecesariamente el procedimiento.
\end{foco}

\subsection{La nulidad del despido por deuda previsional (Ley Bustos) en sede concursal}

Una de las cuestiones más debatidas ha sido la procedencia, en sede concursal, de la
\destacar{nulidad del despido por deuda previsional} del \artcdt{162} ---la llamada Ley
Bustos---, que obliga al empleador con cotizaciones impagas a seguir pagando remuneraciones
hasta convalidar el despido enterando lo adeudado.

La respuesta la da el \destacar{propio texto del \artcdt{163 bis} \Nro 1}: sus reglas se aplican
«de forma preferente a lo establecido en el artículo 162 y en ningún caso se producirá el efecto
establecido en el inciso quinto de dicho artículo». Es decir, la sanción de la Ley Bustos
---y su prolongación hasta la convalidación--- \destacar{no rige} cuando el contrato termina por la
causal concursal del \art{163 bis}: sería incompatible con la naturaleza y la celeridad del
procedimiento, y con el \desasimiento{} que impide seguir devengando remuneraciones tras la
liquidación. El límite no es, pues, una construcción pretoriana, sino una regla legal expresa.

\begin{alerta}[La inaplicabilidad es legal, no jurisprudencial]
Conviene no confundir dos cosas. Que la Ley Bustos \emph{no} opere en el término concursal del
\art{163 bis} es una consecuencia del \emph{texto} de esa norma, no de una línea de fallos. Lo que
la jurisprudencia de unificación ha ido precisando ---en casos de despido \emph{ordinario}, ajenos
al concurso--- son los contornos de la propia Ley Bustos: cuándo se configura y cómo se convalida
(véase enseguida). Invocar el articulado antes que el precedente evita el error de citar como
«límite concursal» fallos que en realidad resolvieron despidos comunes.
\end{alerta}

\subsubsection{La jurisprudencia general de unificación sobre la Ley Bustos}

Al margen del concurso, la Cuarta Sala de la \csup{} ha unificado criterio sobre la Ley Bustos en
pronunciamientos que el operador concursal debe conocer, porque fijan el estándar de la sanción de
la que el \art{163 bis} exceptúa:

\begin{itemize}
  \item \textbf{Compatibilidad con el autodespido} (\rol{40.689-2016}) \parencite{juris-cs-40689-2016}:
  la nulidad del despido es compatible con la acción de despido indirecto o autodespido deducida por
  el trabajador.
  \item \textbf{Cotizaciones sobre horas extraordinarias} (\rol{66.235-2021}, \emph{Rioseco con
  \textsc{efe}}, 26 de diciembre de 2022) \parencite{juris-cs-66235-2021}: el fallo que reconoce
  horas extraordinarias es \destacar{declarativo}, de modo que el impago de cotizaciones sobre ellas
  configura el supuesto de la Ley Bustos.
  \item \textbf{Requisitos copulativos de la convalidación} (\rol{139.538-2022})
  \parencite{juris-cs-139538-2022}: convalidar exige \destacar{dos requisitos copulativos} ---el
  pago efectivo de las cotizaciones adeudadas \emph{y} el envío al trabajador de la carta
  certificada con los comprobantes de pago---; el mero pago, sin la carta, no detiene el devengo de
  las remuneraciones.
\end{itemize}

\begin{foco}[Por qué el concurso corta el reloj y el despido común no]
Fuera del concurso, la sanción de la Ley Bustos corre hasta la convalidación ---con sus dos
requisitos copulativos--- y puede acumular meses de remuneraciones. Dentro del concurso, el
\art{163 bis} la excluye: las cotizaciones impagas anteriores se verifican e integran en el pasivo
de primera clase, pero no generan remuneraciones por nulidad posteriores a la liquidación.
Proyectar la sanción más allá de la resolución ---trasplantando al concurso la regla del despido
común--- infla el crédito y distorsiona el reparto.
\end{foco}

\begin{practica}[Dimensionar el pasivo laboral preferente]
En concursos con planta significativa, el pasivo laboral de primera clase ---remuneraciones,
cotizaciones e indemnizaciones--- puede absorber la mayor parte del activo realizable y, por su
preferencia, pagarse antes que cualquier acreedor valista. Calcular ese pasivo con precisión
---aplicando los topes del \artcc{2472} \Nro 5 y 8 (capítulo \ref{cap:creditos-laborales})--- antes
de proyectar cualquier reparto no es opcional: define cuánto queda para el resto de la masa.
\end{practica}

\subsection{El liquidador en posición de empleador}

El liquidador no sucede al deudor en la calidad de empleador para todos los efectos, pero
asume los deberes de comunicación, emisión de finiquitos y pago. La \superir{} ha
precisado estos deberes mediante instrucciones específicas, sistematizadas en el capítulo
\ref{cap:normativa-superir}.

\subsection{Continuación de las actividades económicas}

Si la junta acuerda la continuidad del giro, la situación laboral se altera: pueden
celebrarse nuevos contratos, cuyos créditos tienen la naturaleza de créditos de la masa y
no de créditos concursales. La distinción determina el orden de pago y debe quedar
documentada desde el primer día.

\begin{alerta}[Créditos de la masa frente a créditos concursales]
Las obligaciones laborales nacidas antes de la resolución de liquidación son créditos
concursales, sujetos a verificación y prelación. Las nacidas después, con ocasión de la
continuidad de giro, son créditos de la masa y se pagan como gasto de administración. La
confusión entre ambas categorías distorsiona íntegramente el reparto.
\end{alerta}

%==============================================================================
\bloquePractico
%==============================================================================

\subsection{Determinación de la nómina laboral}

\begin{checklist}[Antecedentes que el liquidador debe reunir de inmediato]
  \item Contratos de trabajo vigentes y sus anexos.
  \item Libro de remuneraciones de los últimos veinticuatro meses.
  \item Liquidaciones de sueldo del período impago.
  \item Certificado de cotizaciones previsionales y de salud por trabajador.
  \item Certificado de deuda previsional (Previred o equivalente).
  \item Individualización de trabajadores con fuero vigente y su causa.
  \item Contratos con trabajadores a plazo fijo o por obra, con su fecha de término
  pactada.
  \item Registro de asistencia y de feriados utilizados.
  \item Juicios laborales en curso, con tribunal, rol y estado.
\end{checklist}

\subsection{Cuestionario al empleador deudor}

\begin{cuestionario}[Situación laboral de la empresa]
  \pregunta{¿Cuántos trabajadores tiene, con qué antigüedad y bajo qué modalidad
  contractual?}
  {Define la magnitud del pasivo laboral preferente, que se paga antes que cualquier
  acreedor valista.}

  \pregunta{¿Hay remuneraciones o cotizaciones impagas? ¿Desde cuándo?}
  {Las cotizaciones impagas generan responsabilidades autónomas y agravan la posición de
  los administradores.}

  \pregunta{¿Hay trabajadores con fuero maternal, sindical o de otra clase?}
  {Es la principal contingencia del capítulo y debe provisionarse antes de repartir.}

  \pregunta{¿Existe sindicato o instrumento colectivo vigente?}
  {Los instrumentos colectivos pueden contener indemnizaciones convencionales superiores
  a las legales.}

  \pregunta{¿Hay juicios laborales en curso o demandas anunciadas?}
  {Son créditos contingentes que deben reservarse.}

  \pregunta{¿Se pactaron indemnizaciones convencionales por sobre el mínimo legal?}
  {Su tratamiento preferente tiene topes propios (capítulo \ref{cap:creditos-laborales}).}
\end{cuestionario}

\subsection{Cronograma de la desvinculación}

\begin{flujograma}{Secuencia laboral desde la resolución de liquidación}{cap21-secuencia}
  \node[hito] (res) {RESOLUCIÓN DE LIQUIDACIÓN\\[2pt]
    \normalfont Fecha de término de los contratos (\artcdt{163 bis})};
  \node[nodo, below=of res] (nomina)
    {DETERMINACIÓN DE LA NÓMINA LABORAL\\[2pt]
     \normalfont Antigüedad, remuneración, fueros y deuda previsional};
  \node[nodo, below=of nomina] (comunica)
    {COMUNICACIÓN DEL TÉRMINO AL TRABAJADOR Y A LA INSPECCIÓN DEL TRABAJO};
  \node[nodo, below=of comunica] (calculo)
    {CÁLCULO INDIVIDUAL DE INDEMNIZACIONES Y HABERES ADEUDADOS};
  \node[nodo, below=of calculo] (finiquito)
    {EMISIÓN DEL FINIQUITO ELECTRÓNICO\\[2pt]
     \normalfont Plataforma de la \dt{}};
  \node[favorable, below=of finiquito] (pago)
    {PAGO CON PREFERENCIA DE PRIMERA CLASE\\[2pt]
     \normalfont Sujeto al estado de realización del activo};

  \draw[flecha] (res) -- (nomina);
  \draw[flecha] (nomina) -- (comunica);
  \draw[flecha] (comunica) -- (calculo);
  \draw[flecha] (calculo) -- (finiquito);
  \draw[flecha] (finiquito) -- (pago);
\end{flujograma}

\subsection{El deber de comunicación del liquidador y su sanción}

El \artcdt{163 bis} \Nro 1 impone al liquidador comunicar el término de la relación laboral a
cada trabajador ---personalmente o por carta certificada al domicilio señalado en el contrato---
dentro de un plazo \destacar{no superior a seis días hábiles contado desde la fecha de
notificación de la resolución de liquidación} por el tribunal que conoce del procedimiento. A la
comunicación debe adjuntarse un \destacar{certificado de la \superir{}} que indique el inicio del
procedimiento, el tribunal competente, la individualización del proceso y la fecha de la
resolución. Dentro del mismo plazo, el liquidador debe enviar copia de la comunicación a la
respectiva \destacar{Inspección del Trabajo}, que lleva un registro actualizado de las recibidas en
los últimos treinta días hábiles.

La ley precisa que \destacar{el error u omisión en esa comunicación no invalida el término} de la
relación laboral: el contrato termina por el solo ministerio de la ley, y el defecto sólo genera
responsabilidad administrativa. En efecto, la \dt{} constata el incumplimiento ---de oficio o a
petición de parte--- e informa a la \superir{}, que puede sancionar los hechos imputables al
liquidador conforme al \art{338}, sin perjuicio de la responsabilidad penal del Párrafo 7 del
Título IX del Libro Segundo del \cpen{}. Contra la multa procede reposición administrativa dentro
de cinco días (capítulo \ref{cap:institucionalidad}).

\begin{alerta}[Seis días hábiles desde la notificación, no desde el Boletín]
El cómputo no arranca con la publicación en el \boletin{} sino con la \destacar{notificación de la
resolución de liquidación por el tribunal}. Confundir ambos hitos ---que no coinciden--- es la vía
más rápida a una multa de la \superir{}. Es de las primeras actuaciones que deben agendarse el día
de la resolución, no cuando el trabajador reclama.
\end{alerta}

\begin{foco}[El fundamento legal ---no sólo jurisprudencial--- del límite de la Ley Bustos]
El \artcdt{163 bis} \Nro 1 cierra con una declaración decisiva: estas normas
\destacar{«se aplicarán de forma preferente a lo establecido en el artículo 162 y en ningún caso se
producirá el efecto establecido en el inciso quinto de dicho artículo»}. Es decir, la sanción de
nulidad del despido por cotizaciones impagas ---y su prolongación hasta la convalidación--- está
excluida \emph{por texto expreso} en el término concursal. La línea de unificación de la \csup{}
examinada más arriba no innova: aplica lo que la ley ya dispone. Conviene invocar la norma antes
que el precedente.
\end{foco}

%==============================================================================
\bloqueErrores
%==============================================================================

\errorfrecuente{Aplicar el artículo 163 bis a una reorganización}
{Se despide sin causal legal, con la consiguiente demanda por despido injustificado
contra una empresa que intenta reorganizarse.}
{Recordar que la causal opera sólo con la resolución de liquidación.}

\errorfrecuente{Fijar como fecha del término la de la comunicación al trabajador}
{Se altera el cómputo de antigüedad y de haberes, y se genera diferencia reclamable.}
{Usar siempre la fecha de dictación de la resolución de liquidación.}

\errorfrecuente{Repartir el activo sin provisionar las contingencias de fuero}
{Si el trabajador aforado obtiene sentencia favorable, no hay fondos y la responsabilidad
alcanza al liquidador.}
{Reservar el monto de las contingencias antes de cualquier reparto.}

\errorfrecuente{Confundir créditos de la masa con créditos concursales en la continuidad de giro}
{Se altera el orden de pago y se perjudica a los acreedores preferentes.}
{Llevar registro separado desde el primer día de la continuidad.}

\errorfrecuente{Emitir finiquitos en papel}
{Se incumplen las instrucciones vigentes de la \superir{} sobre finiquitos electrónicos.}
{Tramitarlos por la plataforma de la \dt{}, conforme al capítulo
\ref{cap:normativa-superir}.}
